\documentclass{article}
\usepackage[T1]{fontenc}
\usepackage[utf8]{inputenc}
\usepackage[submission]{ismir} % Remove the "submission" option for camera-ready version
\usepackage{amsmath,amssymb,cite,url}
\usepackage{graphicx}
\usepackage{color}

\title{Paper Template for ISMIR \conferenceyear}

% Single address
% To use with only one author or several with the same address
% ---------------
\oneauthor
  {Anonymous Authors}
  {Anonymous Affiliations\\\texttt{anonymous@ismir.net}}

% Two addresses
% --------------
%\twoauthors
%   {First author} {School \\ Department}
%   {Second author} {Company \\ Address}

% Three addresses
% --------------
% \threeauthors
%   {First Author} {Affiliation 1 \\ \texttt{author1@ismir.edu}}
%   {Second Author} {Affiliation 2 \\ \texttt{author2@ismir.edu}}
%   {Third Author} {Affiliation 3 \\ \texttt{author3@ismir.edu}}

% Four or more addresses
% OR alternative format for large number of co-authors
% ------------
% \multauthor
%   {First author$^1$ \hspace{1cm} Second author$^1$ \hspace{1cm} Third author$^2$}
%   {{\bf Fourth author$^3$ \hspace{1cm} Fifth author$^2$ \hspace{1cm} Sixth author$^1$}\\
%   $^1$ Department of Computer Science, University, Country\\
%   $^2$ International Laboratories, City, Country\\
%   $^3$ Company, Address\\
%   {\tt\small CorrespondenceAuthor@ismir.edu, PossibleOtherAuthor@ismir.edu}
%   }

% For the author list in the Creative Common license, please enter author names.
% Please abbreviate the first names of authors and add 'and' between the second to last and last authors.
\def\authorname{F. Author, S. Author, and T. Author}

% Optional: To use hyperref, uncomment the following.
% \usepackage[bookmarks=false,pdfauthor={\authorname},pdfsubject={\pdfsubject},hidelinks]{hyperref}
% Mind the bookmarks=false option; bookmarks are incompatible with ismir.sty.

\sloppy % please retain sloppy command for improved formatting

\begin{document}

\maketitle

\begin{abstract}
Automatic music evaluation is often misaligned with human preference,
especially when intrinsic uncertainty and model confidence are confounded.
We ask a model-centric question:
does a generative music model, through its own internal states,
already know whether the music it produces sounds good to humans?
To answer this, we treat human evaluation scores as ground truth
and probe the model using only its intrinsic signals,
without relying on external audio-content features.
Specifically, we study three complementary signals from a pretrained music model:
two predictive signals at the output level---%
per-step loss and next-token entropy from the logits---%
and one representational signal at the hidden-state level---%
sparse autoencoder (SAE) latent features.
We design controlled settings with single-signal and hybrid variants
to isolate each signal and measure its explanatory power for human ratings.
We further analyze the three signals from complementary angles:
a calibration view on when loss and entropy disagree,
a temporal view on how uncertainty evolves during generation,
and a representational view that links SAE latents
to clips judged good or bad by humans.
Results show that predictive and representational signals
capture complementary aspects of how humans judge generated music,
and that combining them yields a more robust predictor of human evaluation
than any single statistic alone.
\end{abstract}

\section{Introduction}\label{sec:introduction}

Automatic music evaluation still suffers from a persistent gap
between objective scores and human preference.
Existing lines of work have explored score prediction from audio,
uncertainty-aware diagnostics, and representation-based interpretation,
but each direction has distinct strengths and weaknesses
and is often evaluated under different protocols.

This paper takes a model-centric evaluation view.
Instead of using generated audio itself as the primary predictor input,
we use model-internal scoring signals and extracted internal features
to align with human ratings.
Concretely, we study token-level loss dynamics,
next-token entropy from logits, and sparse latent factors from SAE.

Our first contribution is an internal-signal alignment framework:
model-internal signals can be directly aligned with human evaluation targets
across heterogeneous datasets
without requiring an additional audio-content encoder as the main front-end.
Our second contribution is a mechanism-level finding on loss-entropy inconsistency:
low-entropy steps can still produce wrong next-token predictions
when the model is overconfident,
while high-entropy steps indicate diffuse predictive mass and low certainty.
Regions where loss and entropy disagree
are strongly associated with musical surprisal or structural change points.

\section{Background}

Background logic outline for the final cited version:
start from the mismatch between automatic metrics and human judgment in music generation;
review major paradigms (listening tests, audio-input predictors, and model-internal analyses)
and their trade-offs;
summarize prior findings that loss and entropy
are related but not equivalent in sequential prediction;
summarize interpretability literature on sparse latent representations;
then motivate the unresolved gap
of lacking controlled single-signal vs hybrid-signal comparisons under one protocol.

\section{Methodology}

We design a model-aware evaluation pipeline
from five human-evaluation databases to human-aligned prediction targets.
The pipeline includes dataset-specific preprocessing, split construction,
three feature extraction branches (loss, entropy, SAE),
and model training in both single and hybrid settings.

\begin{figure}[t]
  \centering
  \fbox{\parbox{0.92\linewidth}{
  \small
  Five databases $\rightarrow$ preprocessing $\rightarrow$ train/val/test splits $\rightarrow$
  \{loss curve, entropy curve, SAE features\} $\rightarrow$
  \{single models $\times$3, hybrid models $\times$4\} $\rightarrow$
  predicted quality scores $\rightarrow$ correlation with human ratings.
  }}
  \caption{Simplified end-to-end methodology flowchart.}
  \label{fig:pipeline}
\end{figure}

\subsection{Data Processing Across Five Databases}

We use MusicPref, AIME, MusicEval, SongEval, and MusicArena.
MusicEval and SongEval provide continuous ratings
(including multi-dimension or multi-annotator settings),
while MusicPref, AIME, and MusicArena provide preference-style supervision.
We keep each dataset's supervision semantics during preprocessing,
run deterministic train/val/test splits with fixed seeds,
and support both clean and noisy tracks for robustness.

For all pairwise labels, we use one unified objective.
For a pair $(a,b)$, $s_a$ and $s_b$ are latent scalar quality scores,
and $m$ is the ranking margin in the hinge term.
Win and loss use symmetric ReLU penalties,
while tie is optimized by minimizing score distance:
\begin{equation}
\begin{aligned}
\mathcal{L}_{\text{win}} &= \max(0,m+s_b-s_a), \\
\mathcal{L}_{\text{lose}} &= \max(0,m+s_a-s_b), \\
\mathcal{L}_{\text{tie}} &= (s_a-s_b)^2.
\end{aligned}
\end{equation}
After optimization, each latent score $s_i$ is mapped to a common 1--5 range.
Here $\epsilon$ is a small constant for numerical stability:
\begin{equation}
\tilde{s}_i=1+4\cdot\frac{s_i-\min_j s_j}{\max_j s_j-\min_j s_j+\epsilon}.
\end{equation}
For robustness experiments,
noisy labels are generated by bounded Gaussian perturbation,
where $\epsilon_i$ is zero-mean Gaussian noise with variance $\sigma^2$:
\begin{equation}
y_i^{\text{noisy}}=\mathrm{clip}(y_i+\epsilon_i,1,5),\quad \epsilon_i\sim\mathcal{N}(0,\sigma^2).
\end{equation}

\subsection{Feature Extraction}

We extract three complementary feature families.

\textbf{Loss curve.}
Given a token sequence $\{x_t\}_{t=1}^{T}$,
where $t$ is the token-time index and $T$ is sequence length,
teacher-forcing token loss is
\begin{equation}
\begin{aligned}
\ell_t &= -\log p_{\theta}(x_t\mid x_{<t}), \\
\mathbf{l} &= [\ell_1,\ell_2,\ldots,\ell_T].
\end{aligned}
\end{equation}
where $x_{<t}$ is left context
and $p_{\theta}(\cdot\mid x_{<t})$ is the model distribution parameterized by $\theta$.
$\mathbf{l}$ denotes the full loss curve across time.
Besides the raw sequence, we use statistics
such as mean, variance, segment-wise averages, and local peaks.

\textbf{Entropy curve from logits.}
At each step, entropy of the next-token distribution is
\begin{equation}
\begin{aligned}
H_t &= -\sum_{k=1}^{K}p_{\theta}(k\mid x_{<t})\log p_{\theta}(k\mid x_{<t}), \\
\mathbf{h} &= [H_1,H_2,\ldots,H_T].
\end{aligned}
\end{equation}
where $K$ is codebook vocabulary size
and $\mathbf{h}$ is the temporal entropy curve.
To characterize uncertainty-quality interactions,
we also compute clip-level coupling between loss and entropy:
\begin{equation}
\rho_{\ell,H}^{(i)}=\mathrm{corr}(\mathbf{l}^{(i)},\mathbf{h}^{(i)}).
\end{equation}
Here $\rho_{\ell,H}^{(i)}$ is the Pearson correlation within clip $i$.

\textbf{SAE features.}
Let $u_t\in\mathbb{R}^{d}$ be hidden states at time $t$.
SAE produces sparse codes $z_t=f_{\mathrm{enc}}(u_t)$
and reconstructions $\hat{u}_t=f_{\mathrm{dec}}(z_t)$ with objective
\begin{equation}
\mathcal{L}_{\mathrm{SAE}}=
\frac{1}{T}\sum_{t=1}^{T}\|u_t-\hat{u}_t\|_2^2+
\lambda\frac{1}{T}\sum_{t=1}^{T}\|z_t\|_1.
\end{equation}
The first term is reconstruction error,
and $\lambda$ controls sparsity strength via the $\ell_1$ penalty on $z_t$.
From $\{z_t\}$, we build sequence-level and pooled descriptors for downstream prediction.

\subsection{Model Training: Three Single Models and Four Hybrid Models}

We define seven model families.
The three single models are:
\textit{Single-Loss}, \textit{Single-Entropy}, and \textit{Single-SAE}.
The four hybrid models are:
\textit{Loss+Entropy}, \textit{Loss+SAE}, \textit{Entropy+SAE}, and \textit{Loss+Entropy+SAE}.
Fusion is implemented by late concatenation or multi-stream temporal encoding.
In late fusion,
$r_{\text{loss}}$, $r_{\text{ent}}$, and $r_{\text{sae}}$ are branch-level representations
and the final feature is
\begin{equation}
r=[r_{\text{loss}};r_{\text{ent}};r_{\text{sae}}].
\end{equation}

In architecture terms, we use a codebook-aware temporal representation.
For each time step $t$, signals from four codebooks are kept separate and concatenated,
rather than averaged at input time.
With codebook-level feature $f_t^{(c)}$ for codebook index $c\in\{1,2,3,4\}$
and validity mask $m_t\in\{0,1\}$, the model input is
\begin{equation}
v_t=[f_t^{(1)};f_t^{(2)};f_t^{(3)};f_t^{(4)};m_t].
\end{equation}
For hybrid settings,
$f_t^{(c)}$ can include loss, entropy, or SAE-projected channels depending on the branch,
while $m_t$ consistently indicates valid (non-padding) positions.

The mask is applied both during feature preparation and temporal pooling
so that variable-length sequences do not bias the representation.
Here $\phi(\cdot)$ is the temporal encoder and $\epsilon$ avoids division by zero:
\begin{equation}
\tilde{v}_t=m_t v_t,\quad
r=\frac{\sum_{t=1}^{T} m_t\,\phi(\tilde{v}_t)}{\sum_{t=1}^{T} m_t+\epsilon}.
\end{equation}

All predictors are trained as regression to human-aligned targets $y_i\in[1,5]$,
where $\hat{y}_i$ is model output for sample $i$ and $N$ is sample count:
\begin{equation}
\mathcal{L}_{\text{reg}}=\frac{1}{N}\sum_{i=1}^{N}(\hat{y}_i-y_i)^2,
\end{equation}
and optional regularization terms when needed.

\subsection{Evaluation Protocol}

We evaluate predicted scores against human ratings
using Pearson correlation, Spearman correlation, and MSE on held-out sets.
In Pearson, $\bar{\hat{y}}$ and $\bar{y}$ are sample means of predictions and targets:
\begin{equation}
r=\frac{\sum_i(\hat{y}_i-\bar{\hat{y}})(y_i-\bar{y})}
{\sqrt{\sum_i(\hat{y}_i-\bar{\hat{y}})^2}\sqrt{\sum_i(y_i-\bar{y})^2}},
\end{equation}
\begin{equation}
\rho=1-\frac{6\sum_i d_i^2}{n(n^2-1)}.
\end{equation}
In Spearman, $d_i$ is the rank difference for sample $i$
and $n$ is the number of evaluated samples.
All single and hybrid models are compared
under identical splits and preprocessing to ensure fairness.

\section{Results}\label{sec:results}

We report a unified comparison
between large-scale baselines and our internal-signal models
under the same evaluation protocol.
The final numeric results (Pearson/Spearman on clean and noisy tracks)
will be filled after final reruns with fixed seeds and finalized checkpoints.

For aesthetics baselines (CE/CU/PC/PQ),
we apply a score-alignment step before comparison with human MOS.
Let $z_i$ be the raw aesthetics output for sample $i$.
We use an affine mapping to the 1--5 target range,
\begin{equation}
\hat{z}_i=\mathrm{clip}(a z_i+b,1,5),
\end{equation}
where $a$ and $b$ are estimated on the validation split only,
then applied to test predictions.
This keeps baseline and proposed models on the same target scale
while avoiding test-set leakage in calibration.

\begin{table*}[t]
  \centering
  \scriptsize
  \begin{tabular}{|l|l|l|c|c|c|c|}
    \hline
    Group & Model & Input Signals & P (Clean) & S (Clean) & P (Noisy) & S (Noisy) \\
    \hline
    Baseline & Aesthetics-CE & Audio signal (five databases) & TBD & TBD & TBD & TBD \\
    Baseline & Aesthetics-CU & Audio signal (five databases) & TBD & TBD & TBD & TBD \\
    Baseline & Aesthetics-PC & Audio signal (five databases) & TBD & TBD & TBD & TBD \\
    Baseline & Aesthetics-PQ & Audio signal (five databases) & TBD & TBD & TBD & TBD \\
    Baseline & Mean Loss & Mean token loss & TBD & TBD & TBD & TBD \\
    Baseline & Loss Curve CNN & Loss curve + mask & TBD & TBD & TBD & TBD \\
    \hline
    Single & Single-Loss & Loss (4 codebooks) + mask & TBD & TBD & TBD & TBD \\
    Single & Single-Entropy & Entropy (4 codebooks) + mask & TBD & TBD & TBD & TBD \\
    Single & Single-SAE & SAE features + mask & TBD & TBD & TBD & TBD \\
    \hline
    Hybrid & Loss + Entropy & Loss + entropy + mask & TBD & TBD & TBD & TBD \\
    Hybrid & Loss + SAE & Loss + SAE + mask & TBD & TBD & TBD & TBD \\
    Hybrid & Entropy + SAE & Entropy + SAE + mask & TBD & TBD & TBD & TBD \\
    Hybrid & Loss + Entropy + SAE & 3-stream fusion + mask & TBD & TBD & TBD & TBD \\
    \hline
  \end{tabular}
  \caption{%
    Two-column comparison table between baselines and proposed models.
    All pairwise-supervised sources are converted with the same ReLU-based objective
    and aligned to the 1--5 target range before training/evaluation.%
  }
  \label{tab:main_results}
\end{table*}

\section{Analysis}\label{sec:analysis}

While Section~\ref{sec:results} reports aggregate predictive performance,
the three intrinsic signals differ in what they reveal about the model's internal state.
We analyze them from three complementary angles:
a \emph{calibration view} on when loss and entropy disagree,
a \emph{temporal view} on how uncertainty evolves during generation,
and a \emph{representational view} that probes SAE latents
for an interpretable good-versus-bad axis.

\subsection{Calibration View: When Loss and Entropy Disagree}
\label{subsec:calibration}

Loss and entropy are both derived from the same output distribution
$p_\theta(\cdot\mid x_{<t})$ and are often treated as two faces of the same coin.
In a well-calibrated model they co-vary:
high-confidence predictions (low entropy) are usually also correct (low loss).
Modern neural networks, however, are known to be miscalibrated
and systematically overconfident~\cite{guo2017calibration},
and penalizing confident outputs has been shown to act
as an implicit regularizer against this tendency~\cite{pereyra2017confidence}.

We observe the same pattern in generative music models.
Partitioning all token steps into a $2{\times}2$ grid
of \{low, high\} loss $\times$ \{low, high\} entropy
gives four regimes with distinct human-rating signatures:
\begin{itemize}
  \item \textbf{Low loss, low entropy} (confident correct):
    the ideal regime, positively associated with high human ratings.
  \item \textbf{High loss, high entropy} (uncertain wrong):
    the model knows it does not know;
    often co-located with genuine musical surprise
    such as section boundaries, and does not strongly penalize ratings.
  \item \textbf{Low loss, high entropy} (surprise survived):
    diffuse predictive mass where the ground-truth token
    nevertheless receives enough probability.
  \item \textbf{\emph{High loss, low entropy}} (overconfident wrong):
    the mechanism-level failure mode in which the model commits to
    a sharply peaked prediction inconsistent with the ground-truth token.
    This regime is the strongest negative predictor of human rating
    among the four.
\end{itemize}

Concretely, we construct a reliability-style
diagram~\cite{guo2017calibration}
by binning steps according to predictive entropy $H_t$
and reporting the mean token loss $\ell_t$
and mean clip-level human rating within each bin.
A well-calibrated model would produce a monotone increase in $\ell_t$ with $H_t$;
deviations from monotonicity---particularly large $\ell_t$ in the lowest-$H_t$ bins---%
localize the overconfidence regime described above
(see Fig.~\ref{fig:calibration}, TBD).
This analysis gives a principled reason why neither loss nor entropy alone suffices:
each signal misreads one axis of the confidence-correctness plane,
and their \emph{disagreement} is itself an informative feature.

\subsection{Temporal View: Entropy Curves as Uncertainty Signatures}
\label{subsec:temporal}

Aggregate statistics collapse a sequence into a single number
and therefore miss when the model is uncertain.
The temporal shape of $\mathbf{h}=[H_1,\ldots,H_T]$
carries additional information that correlates with human ratings
beyond its mean or variance.

Recent work supports treating entropy as a time-resolved reliability signal.
Entropy has been used at decoding time to suppress distracted
or noisy generations~\cite{qiu2024entropy},
and at the training dynamics level,
an empirical law $R=-a\exp(H)+b$ between policy entropy $H$
and downstream performance $R$ has been established
for reasoning language models~\cite{cui2025entropy},
indicating that entropy and output quality are structurally,
not coincidentally, linked.
In the generative setting, teacher-forced training
and free-running inference produce systematically different
sequence distributions~\cite{wiseman2016sequence},
meaning that the teacher-forced loss curve alone
can misrepresent the model's behavior under its own rollouts.
Entropy, computed directly from the predictive distribution,
is comparatively less affected by this exposure-bias gap
and thus a useful complementary temporal signal.

Empirically, we analyze entropy trajectories in three ways:
\begin{enumerate}
  \item \textbf{Rating-stratified averages.}
    Clips are grouped by human-rating quartile
    and mean entropy curves are plotted after length normalization.
    Higher-rated clips exhibit entropy trajectories with moderate local variation
    around structural boundaries;
    lower-rated clips tend to show either near-flat entropy
    (symptomatic of mode collapse or repetition)
    or erratic oscillation.
  \item \textbf{Curve statistics as features.}
    We extract temporal descriptors---variance, lag-$k$ autocorrelation,
    peak density, and spectral centroid of the entropy curve---%
    and ablate them against the raw curve.
    The lift they provide over $\{\mathrm{mean}(H),\mathrm{var}(H)\}$
    quantifies how much \emph{temporal structure}, not level, drives the prediction.
  \item \textbf{Loss--entropy coupling dynamics.}
    Using the clip-level coupling $\rho_{\ell,H}^{(i)}$ introduced earlier,
    we inspect how $\rho$ distributes across rating quartiles
    and where within the sequence the two curves diverge the most.
\end{enumerate}
Together, these diagnostics convert entropy from a scalar summary
into a time-resolved uncertainty signature
that reflects how the model navigates musical structure during generation.

\subsection{Representational View: SAE Latents as Good-vs-Bad Probes}
\label{subsec:sae}

Loss and entropy are logit-level summaries:
they compress the full predictive distribution into one number per step.
SAE features instead act on the hidden state $u_t$
and provide a \emph{representational} readout
that is not available from the output alone.
This motivates a third analysis question:
do the sparse latents $\{z_t^{(k)}\}$ contain directions
that already separate human-preferred from human-dispreferred clips,
without any supervision from human ratings during SAE training?

To answer this, we perform a latent-level correlation analysis.
For each SAE unit $k$, let
\begin{equation}
  \bar{z}^{(k,i)}=\frac{1}{T_i}\sum_{t=1}^{T_i} z_t^{(k,i)}
\end{equation}
denote its mean activation on clip $i$, and let $y_i$ be the clip's human rating.
We score each unit by its Pearson correlation
$r_k=\mathrm{corr}(\{\bar{z}^{(k,i)}\}_i,\{y_i\}_i)$
and rank units by $|r_k|$
to isolate the most rating-aligned directions,
separately for $r_k>0$ and $r_k<0$.

Three probes follow:
\begin{itemize}
  \item \textbf{Top-ranked latent profiles.}
    For the top positive and top negative latents,
    we report activation distributions
    on clips in the highest and lowest rating quartiles.
    Clear bimodality indicates that the pretrained model
    already encodes a good-vs-bad axis that can be read out linearly.
  \item \textbf{Activation-ranked retrieval.}
    For a given top latent $k^\star$,
    we retrieve the clips with the largest $\bar{z}^{(k^\star,i)}$
    and inspect them qualitatively.
    This turns an abstract latent index
    into an interpretable musical property
    (e.g., stable harmonic motion vs.\ abrupt timbre changes).
  \item \textbf{Latent-subset ablation.}
    We compare Single-SAE trained on the full SAE dictionary
    against Single-SAE restricted to the top-$K$ rating-correlated latents.
    The relative gap quantifies how concentrated
    the human-aligned information is within the sparse dictionary.
\end{itemize}

Unlike the calibration and temporal views---%
which explain \emph{when} the model is confident or surprised---%
the representational view asks a different question:
\emph{whether the model, through its hidden states alone,
already has an internal notion of what humans will rate as good or bad music.}
Our results suggest that this notion exists and is at least partially linear,
which is why SAE features yield gains that are complementary to,
rather than overlapping with, loss and entropy.

\section{Future Work}

\section{Acknowledgement}

To be completed for the camera-ready version (omit in anonymous submission).

\bibliographystyle{IEEEtran}
\bibliography{ISMIRtemplate}

\end{document}
